\chapter{微积分基本定理}
\label{chap:fundamental-theorem-calculus}

Riemann 积分由函数值的有限和定义，导数由局部差商定义。两者的联系不是
定义的一部分，需要定理来建立。积分上限函数把全区间上的累积量重新变成
一点附近可比较的函数；其增量恰是一个短区间上的积分，这使差商与局部平均
联系起来。

\section{积分上限函数}

设 \(f\) 在 \([a,b]\) 上 Riemann 可积，定义
\[
 F(x)=\int_a^x f(t)\,dt,\qquad x\in[a,b],
\]
其中 \(x<a\) 的反向积分约定不会在本节出现。

\begin{theorem}[积分上限函数的连续性]
\label{thm:indefinite-integral-lipschitz}
若 \(\abs f\le M\)，则
\[
 \abs{F(y)-F(x)}\le M\abs{y-x}
 \qquad(x,y\in[a,b]).
\]
因此 \(F\) 是 Lipschitz 连续函数。
\end{theorem}

\begin{proof}
不妨设 \(x<y\)。区间可加性和积分估计给出
\[
 F(y)-F(x)=\int_x^y f(t)\,dt,
\qquad
 \abs{F(y)-F(x)}\le M(y-x).
\]
\end{proof}

这个结论只需 \(f\) 可积，不需连续。可导性则取决于基点附近的局部行为。

\begin{theorem}[连续点处的求导]
\label{thm:indefinite-integral-derivative-continuity-point}
若 \(f\) Riemann 可积并在 \(x_0\in(a,b)\) 连续，则
\[
 F'(x_0)=f(x_0).
\]
端点连续时，相应内侧导数也成立。
\end{theorem}

\begin{proof}
当 \(h\ne0\) 且 \(x_0+h\in[a,b]\) 时，
\[
 \frac{F(x_0+h)-F(x_0)}h-f(x_0)
 =\frac1h\int_{x_0}^{x_0+h}\bigl(f(t)-f(x_0)\bigr)\,dt.
\]
由积分估计，
\[
 \abs{\frac{F(x_0+h)-F(x_0)}h-f(x_0)}
 \le
 \sup_{\min(x_0,x_0+h)\le t\le\max(x_0,x_0+h)}
 \abs{f(t)-f(x_0)}.
\]
右端由连续性趋于零。
\end{proof}

\begin{counterexample}
令 \(f=\mathbf1_{[0,1]}\) 于 \([-1,1]\)。则
\[
 F(x)=\int_{-1}^x f(t)\,dt
 =
 \begin{cases}
 0,&x\le0,\\
 x,&x\ge0.
 \end{cases}
\]
\(f\) 可积，但 \(F\) 在 \(0\) 不可导。积分上限函数在被积函数的连续点
可导，并不保证在每个点可导。
\end{counterexample}

\section{微积分基本定理与 Newton--Leibniz 公式}

\begin{theorem}[微积分基本定理第一部分]
\label{thm:ftc-first}
若 \(f\in C([a,b])\)，则
\[
 F(x)=\int_a^x f(t)\,dt
\]
在 \((a,b)\) 上可导且 \(F'=f\)，并在端点具有相应单侧导数。
\end{theorem}

这一定理同时给出连续函数原函数的存在性。它的证明已包含在
\theoremref{thm:indefinite-integral-derivative-continuity-point} 中。

\begin{theorem}[Newton--Leibniz 公式]
\label{thm:newton-leibniz}
设 \(G\) 在 \([a,b]\) 上连续、在 \((a,b)\) 上可导，并且
\(G'\) 在 \([a,b]\) 上有一个 Riemann 可积的延拓 \(g\)。则
\[
 \int_a^b g(x)\,dx=G(b)-G(a).
\]
特别地，若 \(G\in C^1([a,b])\)，则
\[
 \int_a^b G'(x)\,dx=G(b)-G(a).
\]
\end{theorem}

\begin{proof}
给定分割 \(P:a=x_0<\cdots<x_n=b\)。由 Lagrange 中值定理，对每个
\(i\) 存在 \(\xi_i\in(x_{i-1},x_i)\)，使
\[
 G(x_i)-G(x_{i-1})=G'(\xi_i)\Delta x_i
 =g(\xi_i)\Delta x_i.
\]
求和得到
\[
 G(b)-G(a)=S(g,\dot P).
\]
令分割网格趋于零，右端趋于 \(\int_a^b g\)。
\end{proof}

\begin{corollary}[微积分基本定理第二部分]
若 \(f\in C([a,b])\)，而 \(G\) 是 \(f\) 的任一原函数，则
\[
 \int_a^b f(x)\,dx=G(b)-G(a).
\]
\end{corollary}

“微分与积分互为逆运算”应按两个方向分别理解：连续函数积分后再求导，
恢复原函数；具有 Riemann 可积导数的函数求导后再积分，恢复端点增量。
两个方向的假设不同。

\section{原函数存在性的条件}

连续性是原函数存在的充分条件，不是必要条件。例如
\[
 F(x)=
 \begin{cases}
 x^2\sin(1/x^2),&x\ne0,\\
 0,&x=0
 \end{cases}
\]
处处可导，而 \(F'\) 在 \(0\) 不连续。另一方面，Riemann 可积性并不足以
保证原函数存在。

\begin{proposition}
\label{prop:integrable-no-antiderivative}
定义
\[
 f(x)=
 \begin{cases}
 0,&x<0,\\
 1,&x\ge0
 \end{cases}
\qquad(-1\le x\le1).
\]
则 \(f\) Riemann 可积，但不存在区间 \((-1,1)\) 上的可导函数
\(G\) 满足 \(G'=f\)。
\end{proposition}

\begin{proof}
可积性来自单调函数定理。若存在 \(G\)，则导函数 \(G'\) 必有 Darboux
性质；但 \(f\) 在 \(0\) 处从 \(0\) 跳到 \(1\)，没有介值性，矛盾。
\end{proof}

\begin{proposition}[原函数的唯一性]
若 \(F,G\) 是同一区间上函数 \(f\) 的原函数，则 \(F-G\) 为常函数。
\end{proposition}

\begin{proof}
\((F-G)'=0\)，由中值定理的零导数判据即得。
\end{proof}

积分上限函数提供一种折中的结论：对任意 Riemann 可积 \(f\)，它始终连续，
且在 \(f\) 的每个连续点可导；只有当 \(f\) 处处连续时，才立即得到处处
原函数。

\section{积分中值定理}

\begin{theorem}[第一积分中值定理]
\label{thm:first-integral-mean-value}
若 \(f\in C([a,b])\)，则存在 \(\xi\in[a,b]\)，使
\[
 \int_a^b f(x)\,dx=f(\xi)(b-a).
\]
\end{theorem}

\begin{proof}
令 \(m=\min f\)、\(M=\max f\)。积分的序性质给出
\[
 m(b-a)\le\int_a^b f\le M(b-a).
\]
故平均值
\[
 A=\frac1{b-a}\int_a^b f
\]
介于 \(m,M\) 之间。由连续函数介值定理，存在 \(\xi\) 使
\(f(\xi)=A\)。
\end{proof}

\begin{theorem}[带权积分中值定理]
\label{thm:weighted-integral-mean-value}
若 \(f\in C([a,b])\)，\(g\) Riemann 可积且 \(g\ge0\)，则存在
\(\xi\in[a,b]\)，使
\[
 \int_a^b f(x)g(x)\,dx
 =f(\xi)\int_a^b g(x)\,dx.
\]
\end{theorem}

\begin{proof}
若 \(\int g=0\)，因 \(g\ge0\)，由序估计可得 \(\int fg=0\)，任取
\(\xi\) 即可。若 \(\int g>0\)，由
\[
 m g\le fg\le M g
\]
积分并除以 \(\int g\)，所得加权平均仍介于 \(m,M\) 之间，再用介值定理。
\end{proof}

连续性不能简单删去。若 \(f\) 只取 \(0,1\) 两值，而它的积分平均值严格
介于二者之间，就不存在这样的 \(\xi\)。

\section{分部积分和变量代换}

\begin{theorem}[分部积分]
\label{thm:riemann-integration-by-parts}
若 \(u,v\in C^1([a,b])\)，则
\[
 \int_a^b u(x)v'(x)\,dx
 =u(b)v(b)-u(a)v(a)-\int_a^b u'(x)v(x)\,dx.
\]
\end{theorem}

\begin{proof}
乘积法则给出
\((uv)'=u'v+uv'\)。对两边应用 Newton--Leibniz 公式并移项。
\end{proof}

\begin{theorem}[变量代换]
\label{thm:riemann-substitution}
设 \(\phi\in C^1([\alpha,\beta])\)，\(f\) 在包含
\(\phi([\alpha,\beta])\) 的区间上连续，则
\[
 \int_\alpha^\beta f(\phi(t))\phi'(t)\,dt
 =\int_{\phi(\alpha)}^{\phi(\beta)}f(x)\,dx.
\]
\end{theorem}

\begin{proof}
取 \(f\) 的原函数 \(F\)。链式法则给出
\((F\circ\phi)'=(f\circ\phi)\phi'\)。Newton--Leibniz 公式于是得到
\[
 \int_\alpha^\beta(f\circ\phi)\phi'
 =F(\phi(\beta))-F(\phi(\alpha))
 =\int_{\phi(\alpha)}^{\phi(\beta)}f.
\]
\end{proof}

这里不要求 \(\phi\) 单调或一一对应，因为左右两端都是带方向的积分，
曲线折返造成的贡献由 \(\phi'\) 的符号自动抵消。若把右端改成普通“面积”
或擅自把 \(\phi'\) 换成 \(\abs{\phi'}\)，结论就不同。

\section{Taylor 公式的积分余项与积分 Jensen 不等式}

\begin{theorem}[Taylor 公式的积分余项]
\label{thm:taylor-integral-remainder}
若 \(f\in C^{n+1}\) 于包含 \(a,x\) 的闭区间，则
\[
 f(x)=\sum_{k=0}^n\frac{f^{(k)}(a)}{k!}(x-a)^k
 +\frac1{n!}\int_a^x f^{(n+1)}(t)(x-t)^n\,dt.
\]
\end{theorem}

\begin{proof}
只写 \(a<x\)；反向积分情形相同。对积分连续使用分部积分：
\[
 \frac1{n!}\int_a^x f^{(n+1)}(t)(x-t)^n\,dt
 =-\frac{f^{(n)}(a)}{n!}(x-a)^n
  +\frac1{(n-1)!}\int_a^x f^{(n)}(t)(x-t)^{n-1}\,dt.
\]
重复 \(n\) 次，最后一项为
\(\int_a^x f'(t)\,dt=f(x)-f(a)\)。整理即得公式。这也补全了
\chapterref{chap:taylor-theory} 中式
\eqref{eq:taylor-integral-remainder} 的证明。
\end{proof}

\begin{corollary}
若 \(\abs{f^{(n+1)}}\le M\) 于连接 \(a,x\) 的区间，则
\[
 \abs{R_{n,a}f(x)}
 \le \frac{M}{(n+1)!}\abs{x-a}^{n+1}.
\]
\end{corollary}

\begin{proof}
对积分余项取绝对值。令
\[
 I_{a,x}=[\min\{a,x\},\max\{a,x\}],
\]
则定向积分的绝对值不超过在 \(I_{a,x}\) 上对绝对值的积分，因而
\[
 \frac{M}{n!}\int_{I_{a,x}}\abs{x-t}^n\,dt
 =\frac{M}{(n+1)!}\abs{x-a}^{n+1}.
\]
\end{proof}

\begin{theorem}[积分 Jensen 不等式]
\label{thm:integral-jensen-riemann}
设 \(\Phi\) 在区间 \(J\) 上连续且凸，\(f:[a,b]\to J\) 连续，
\(w:[a,b]\to[0,+\infty)\) Riemann 可积且
\(\int_a^b w>0\)。若 \(\Phi\circ f\) 可积，则
\[
 \Phi\!\left(
 \frac{\int_a^b fw}{\int_a^b w}
 \right)
 \le
 \frac{\int_a^b \Phi(f)w}{\int_a^b w}.
\]
\end{theorem}

\begin{proof}
对细标记分割，把非负数
\[
 \lambda_i=
 \frac{w(\xi_i)\Delta x_i}
      {\sum_j w(\xi_j)\Delta x_j}
\]
作为权重，有限 Jensen 不等式给出相应有限和不等式。由于
\(\int w>0\)，分母对充分细分割保持正。令网格趋于零，Riemann 和分别
趋于三个积分。有限和中的加权平均趋于
\((\int fw)/(\int w)\)，而 \(\Phi\) 的连续性允许在左端通过极限，
得到结论。
\end{proof}

取等权函数和线性参数，可以从凸性得到区间中点、积分平均与端点值之间的
另一组基本估计。

\begin{theorem}[Hermite--Hadamard 不等式]
\label{thm:hermite-hadamard}
设 \(f:[a,b]\to\R\) 连续且凸，其中 \(a<b\)。则
\begin{equation}
 f\!\left(\frac{a+b}{2}\right)
 \le \frac1{b-a}\int_a^b f(x)\,dx
 \le \frac{f(a)+f(b)}2.
 \label{eq:hermite-hadamard}
\end{equation}
若 \(f\) 严格凸，则两侧不等式都严格。
\end{theorem}

\begin{proof}
令 \(m=(a+b)/2\)。对 \(x\in[a,m]\)，点 \(m\) 是 \(x\) 与
\(a+b-x\) 的中点，故
\[
 2f(m)\le f(x)+f(a+b-x).
\]
在 \([a,m]\) 上积分，并对第二项作变量代换，得到
\[
 (b-a)f(m)\le\int_a^b f(x)\,dx.
\]
另一方面，把 \(x\in[a,b]\) 写成
\(x=(1-t)a+tb\)，其中 \(t=(x-a)/(b-a)\)。凸性给出
\[
 f(x)\le\frac{b-x}{b-a}f(a)+\frac{x-a}{b-a}f(b).
\]
两边在 \([a,b]\) 上积分即得右侧不等式。若 \(f\) 严格凸，则上述点对
在除端点或中点外均为真凸组合，严格不等式在正长度的区间上成立，积分后
仍严格。
\end{proof}

\section*{习题}
\addcontentsline{toc}{section}{习题}

\subsection*{A组　积分上限函数与基本定理}
\begin{enumerate}[label=\textbf{\thechapter.\arabic*.}]
 \exerciseitem{ftc-indefinite-step-function} 对本章阶梯函数反例，逐点求积分
 上限函数的左右导数，并与被积函数的连续点比较。
 \exerciseitem{ftc-continuity-point-general} 设 \(f\) 可积，证明积分上限函数
 在 \(f\) 的任一连续点可导；要求同时处理 \(h>0\) 与 \(h<0\)。
 \exerciseitem{ftc-antiderivative-difference} 用中值定理证明同一函数的任意
 两个原函数只差常数，并说明定义域若不是区间时结论应怎样修改。
 \exerciseitem{ftc-integral-zero-continuous} 若连续函数 \(f\) 满足
 \(\int_a^x f(t)\,dt=0\) 对所有 \(x\in[a,b]\) 成立，证明 \(f\equiv0\)。
 \exerciseitem{ftc-derivative-integrable-newton} 设 \(F\) 可导且 \(F'\) Riemann
 可积。直接用每段中值点构成的 Riemann 和证明 Newton--Leibniz 公式。
 \exerciseitem{ftc-integrable-no-primitive-darboux} 构造一个只有一个跳跃点的
 Riemann 可积函数，并用 Darboux 定理证明它没有原函数。
\end{enumerate}

\subsection*{B组　积分恒等式与估计}
\begin{enumerate}[label=\textbf{\thechapter.\arabic*.},resume]
 \exerciseitem{integral-mean-value-weight-sign} 举例说明带权积分中值定理中
 \(g\ge0\) 不能简单删去。
 \exerciseitem{integral-average-between-extrema} 证明连续函数的积分平均值严格
 介于最小值和最大值之间，当且仅当函数不是常函数。
 \exerciseitem{hermite-hadamard} 不引用正文结论，分别从中点凸性和弦下估计
 证明 Hermite--Hadamard 不等式；再证明：若 \(f\) 严格凸且 \(a<b\)，
 则两个不等号都严格。
 \exerciseitem{integration-by-parts-log} 用分部积分计算
 \(\int_1^b\log x\,dx\)，并只引用已经证明的定理说明每一步合法。
 \exerciseitem{substitution-nonmonotone} 取 \(\phi(t)=t^2\) 于 \([-1,1]\)，
 分别对 \(f(x)=1\) 与 \(f(x)=x\) 核对变量代换公式，解释折返抵消。
 \exerciseitem{substitution-oriented-integral} 从变量代换公式推出
 \[
  \int_a^b f(x)\,dx=-\int_{-b}^{-a}f(-u)\,du,
 \]
 并导出奇偶函数在对称区间上的积分公式。
 \exerciseitem{ftc-taylor-integral-remainder-bound} 从积分余项推出 Lagrange 型
 一致误差界，并说明该推导只给出估计，不自动给出某个中间点公式。
 \exerciseitem{integral-jensen-square} 以 \(\Phi(u)=u^2\) 推出
 \[
 \left(\frac1{b-a}\int_a^b f\right)^2
 \le\frac1{b-a}\int_a^b f^2,
 \]
 并判定连续函数情形的等号条件。
\end{enumerate}

\subsection*{C组　条件与推广}
\begin{enumerate}[label=\textbf{\thechapter.\arabic*.},resume]
 \exerciseitem{ftc-absolute-value-antiderivative} 证明 \(\abs x\) 有
 \(C^1\) 原函数，但符号函数没有原函数；比较两个结论。
 \exerciseitem{ftc-integral-upper-function-convex} 若 \(f\) 连续且单调不减，
 证明 \(F(x)=\int_a^x f\) 为凸函数；给出严格凸的充分条件。
 \exerciseitem{ftc-second-derivative-integral-reconstruction} 设
 \(f\in C^2([a,b])\)。证明
 \[
 f(x)=f(a)+f'(a)(x-a)+\int_a^x(x-t)f''(t)\,dt
 \]
 并由此重新证明 \(f''\ge0\Rightarrow f\) 凸。
 \exerciseitem{integral-jensen-weighted-equality} 若 \(\Phi\) 严格凸、
 \(f,w\) 连续、\(w>0\)，证明积分 Jensen 不等式取等当且仅当
 \(f\) 为常函数。
 \optionalexerciseitem{ftc-derivative-discontinuity-primitive} 研究
 \(F(x)=x^2\sin(1/x^2)\)（在 \(0\) 处补零）的导函数，证明该导函数虽不连续
 却有原函数；再说明为什么它的导数在包含 \(0\) 的区间上不是 Riemann 可积。
\end{enumerate}
